% Stage 8 supplement: learning, implementation, and verification layer.
% This is additive: it audits and connects the preceding modules.
\chapter{Learning and Verification Layer}
\label{ch:stage8-learning-verification}

This chapter turns the preceding supplements into a student-facing course
path.  It does not replace any topic.  Instead, it supplies worked examples,
proof and implementation tasks, source-course tags, research directions, and
an assessment rubric.  A complete solution is expected to state the input
model, random choices, invariant, guarantee, error probability, runtime, and
space usage.

\section{Worked example: amplifying a randomized decision algorithm}

Suppose an algorithm returns the correct answer with probability at least
$2/3$, independently on each run.  Run it $r$ times and return the majority.
If each run is a Bernoulli trial with success probability $p\geq2/3$, then
Hoeffding's inequality gives
\[
 \Pr[\text{majority is wrong}]
 \leq \exp\left(-2r(p-1/2)^2\right)\leq e^{-r/18}.
\]
Thus $r\geq18\ln(1/\delta)$ gives failure probability at most $\delta$.
The independence assumption is part of the theorem: reusing the same random
seed or correlated sketches does not justify this calculation.  For a
one-sided Monte Carlo algorithm, repetition can instead return the best
candidate or the logical OR/AND appropriate to its one-sided error.

\paragraph{Student check.}
Write down which output rule is valid for (i) a verifier with false positives,
(ii) a verifier with false negatives, and (iii) a two-sided estimator.  Give a
counterexample to blindly taking a majority in case (iii).

\section{Worked example: a streaming heavy hitter}

Let $f_i$ be the frequency of item $i$ in a stream of length $m$.  A Count-Min
Sketch with $d$ independent hash rows and width $w$ stores counters.  Query
item $i$ by the minimum counter among its $d$ hashed locations.  Since each
collision contribution is nonnegative, the estimate never underestimates.
For a fixed query, choosing $w\geq e/\varepsilon$ makes a row's collision
noise exceed $\varepsilon m$ with probability at most $1/e$ by Markov's
inequality.  Choosing $d\geq\ln(1/\delta)$ makes the minimum exceed that
threshold with probability at most $\delta$.

The contract is therefore an additive bound
\[
 f_i\leq\widehat f_i\leq f_i+\varepsilon m
\]
for one fixed query, using $O(\varepsilon^{-1}\log(1/\delta))$ counters.
To guarantee all queries simultaneously, the hash width or failure budget
must be adjusted with a union bound.  This distinction is a common source of
incorrect student claims.

\section{Worked example: sampling, burn-in, and effective samples}

For a Markov chain with stationary distribution $\pi$, let $X_t$ be the state
after burn-in and let $g$ be bounded in $[0,1]$.  If the total-variation
distance of $X_t$ from $\pi$ is at most $\eta$, then
\[
 |\mathbb E[g(X_t)]-\mathbb E_\pi[g]|\leq\eta.
\]
This is a bias statement, not a variance statement.  Consecutive samples are
correlated; collecting $N$ states does not automatically give variance
$1/(4N)$.  A valid report must state a mixing or autocorrelation bound, a
burn-in length, and a thinning or batch-means method if an error bar is
claimed.  The same checklist applies to Metropolis sampling, Glauber dynamics,
randomized 2-SAT, and graph-coloring chains.

\section{Proof laboratory}

The following exercises form a proof progression.

\begin{problem}[Second moment]
Let $X$ count objects surviving a random construction.  Prove the
Paley--Zygmund lower bound
\[
 \Pr[X\geq\theta\mathbb EX]\geq
 (1-\theta)^2\frac{(\mathbb EX)^2}{\mathbb E[X^2]}
\]
for $0\leq\theta<1$, and apply it to a random graph property.
\end{problem}

\begin{problem}[Coupling]
Construct a coupling for two lazy random walks on the hypercube and use the
expected Hamming distance to derive an explicit total-variation upper bound.
Identify the step at which a union bound is used.
\end{problem}

\begin{problem}[Derandomization]
Given a randomized algorithm whose conditional expected cost can be computed
after each exposed bit, prove that choosing the next bit with no larger
conditional expectation yields a deterministic outcome no worse than the
initial expectation.
\end{problem}

\begin{problem}[Geometry]
Use backward analysis to bound the expected number of conflict updates in a
randomized incremental construction when each partial structure has at most
$b$ boundary features.  State which assumptions fail in the presence of
degenerate input and how symbolic perturbation repairs them.
\end{problem}

\section{Implementation laboratory}

Each task should include a deterministic seed option, input validation, a
counter of random draws, and a test comparing empirical behavior with the
proved bound.

\begin{enumerate}
\item Implement amplification for a noisy decision oracle and plot failure
      probability against the number of repetitions.
\item Implement Count-Min Sketch and verify the one-sided additive guarantee
      on adversarial collision-heavy streams.
\item Implement Wilson's uniform-spanning-tree sampler and compare edge
      frequencies with effective-resistance predictions.
\item Implement a lazy Metropolis chain on a small state space; estimate
      burn-in using total variation against an exactly computed stationary law.
\item Implement Karger's contraction algorithm with multiedge support and
      compare independent repetition with recursive contraction.
\item Implement randomized distance-to-subset coordinates and report observed
      embedding distortion, dimension, and construction time.
\end{enumerate}

\section{Research extensions}

The following are deliberately open-ended extensions rather than missing
prerequisites: limited-independence generators for streaming, spectral
independence for spin systems, concentration for dependent random variables,
near-linear-time Laplacian solvers, variance-optimal rare-event proposals,
oblivious routing, differential privacy, and randomized numerical linear
algebra.  A research report should distinguish a theorem proved in the book,
a theorem cited from the literature, and an empirical conjecture.

\section{Course-source tags and reading map}

The tags below identify the expected course family, not exclusive ownership.
\begin{center}
\begin{tabular}{p{0.27\textwidth}p{0.25\textwidth}p{0.38\textwidth}}
\textbf{Module} & \textbf{Tags} & \textbf{Student outcome}\\
Probability and concentration & MIT / Stanford / Berkeley & derive bounds and audit independence\\
Classical algorithms & MIT / Stanford & design, amplify, and analyze randomized algorithms\\
Hashing and streaming & Stanford / Berkeley & give space, error, and failure guarantees\\
Graphs and derandomization & MIT / Berkeley & use probabilistic existence and constructive choices\\
Markov chains and sampling & MIT / Stanford / Berkeley & separate stationarity, mixing, bias, and variance\\
Geometry and counting & Stanford / Berkeley & connect sampling, embeddings, and estimators\\
\end{tabular}
\end{center}
The course map is a guide for review: MIT emphasizes probabilistic methods,
Markov chains, and advanced algorithmic analysis; Stanford emphasizes the
randomized-algorithm toolkit and applications; Berkeley emphasizes probability,
graph processes, and algorithmic problem solving.  The book preserves the
overlap as a coherent path rather than splitting it into incompatible tracks.

\section{Assessment rubric}

A satisfactory solution earns credit only when all of the following are
visible: (1) a precise random experiment; (2) a correctness argument; (3) a
quantitative probability or approximation guarantee; (4) runtime and memory
complexity; (5) handling of edge cases and degeneracies; and (6) a reproducible
experiment or test when implementation is requested.  A polished capstone
should compare two methods on the same inputs and explain where the theorem's
assumptions do and do not match the experiment.

\section{Final verification checklist}

Before considering the book complete, a student or instructor should verify:
\begin{enumerate}
\item every chapter compiles from its LaTeX source;
\item every new supplement is included exactly once in the full book;
\item each algorithm states its model, guarantee, error, runtime, and space;
\item exercises include probability, proof, implementation, and extension work;
\item experiments use reproducible seeds and report uncertainty;
\item no existing chapter or topic was removed during expansion.
\end{enumerate}

\paragraph{Completion criterion.}
The expansion is educationally complete when a student can move from a
problem statement to a random experiment, proof, implementation, and
empirical audit, while knowing precisely which claim is theoretical and which
is experimental.
